\chapter{Álgebra}

Aquí exploraremos los conceptos básicos
relacionados al álgebra. El álgebra
aparece en todas partes de la matemática,
y es importante conocer lo básico
ya que en particular se puede aplicar a
infinidad de problemas.


\section{Identidades}
Manejar expresiones algebraicas puede ser
bastante complicado si no se es capaz
de simplificar la misma.

E aquí algunos métodos y sustituciones
útiles.

\subsection{Trinomio cuadrado perfecto}
El trinomio cuadrado perfecto es quizás
el caso de factoreo más común de todos,
y tan solo consiste en saber que
\[(a + b)^2 = a^2 + 2ab + b^2.\]

En los problemas, a veces es conveniente
expandir la expresión $(a + b)^2$
y en otros casos resulta útil simplificar
$a^2 + 2ab + b^2$ como $(a + b)^2$.

\subsection{Super identidad}
Un amigo me contó que esta simple
transformación aparece en muchos lugares
inesperados y especialmente en teoría de números:
\[ab + cd = a(b - c) + c(d + a)\]

\subsection{Identidad de Sophie Germain}
La identidad de Sophie Germain establece lo siguiente:
\[a^4 + 4b^4 = (a^2 + 2b^2 + 2ab)(a^2 + 2b^2 - 2ab)\]

\section{Problemas}
\begin{problem}
Muestre que
\[1^2 + 3^2 + \cdots + n^2 = \frac{n(n + 1)(2n + 1)}6.\]
\end{problem}

\begin{problem}
Demuestre que la suma de los primeros
$n$ números impares es igual a $n^2$.
\end{problem}

\begin{problem}
Resuelva el sistema de ecuaciones:
\begin{equation*}
	\begin{cases}
		a^2 + b^2 + ab = 7 \\
		a - b          = 2.
	\end{cases}
\end{equation*}
\end{problem}

\begin{problem}
Si $a$ es un número real distinto de $0$, encuentre el valor de la expresión:
\[\frac{a + a^2 + a^3 + \cdots + a^n}{\frac{1}{a} + \frac{1}{a^2} + \frac{1}{a^3} + \cdots + \frac{1}{a^n}}\]
\end{problem}

\begin{problem}[Omapa 2023 - Ronda Departamental]
Encuentre el valor de la expresión
\[\frac1{\sqrt1 + \sqrt2} + \frac1{\sqrt2 + \sqrt3} + \cdots + \frac1{\sqrt{2022} + \sqrt{2023}}.\]
\end{problem}

\begin{problem}[Libro ruso]
Demuestre que si los números positivos $a$, $b$ y $c$
forman una progresión aritmética,
entonces los números
\[\frac1{\sqrt b + \sqrt c}, \frac1{\sqrt c + \sqrt a}, \frac1{\sqrt a + \sqrt b}\]
también forma una progresión aritmética.

\end{problem}

\begin{problem}
Encuentre el valor de la expresión:
\[\frac1{1\cdot2} + \frac1{2\cdot3} + \frac1{3\cdot4} + \cdots + \frac1{n(n + 1)}.\]
\addhints{telescopica1}.
\end{problem}
\section{Ecuaciones}
